\ifdefined\gkver\else\def\gkver{student}\fi
\documentclass[\gkver]{gaokaozhenti}
\begin{document}

\chapter{2012年山东理综}
%% paperType: 理综物理部分

\begin{enumerate}
\item
%% number: 14
%% paperName: 2012年山东理综第 14 题 5 分
%% typeId: 2
%% type: 多选题
%% chapter: 其他
%% point: 物理学史
%% method: 无
%% score: 5
%% degree: 650
%% duplicateId: 0
%% body:
以下叙述正确的是\xzanswer{AD}
\fourchoices[answer=AD]
{法拉第发现了电磁感应现象}
{惯性是物体的固有属性，速度大的物体惯性一定大}
{牛顿最早通过理想斜面实验得出力不是维持物体运动的必然结果}
{感应电流遵从楞次定律所描述的方向，这是能量守恒定律的必然结果}
%% answer: AD
%% memo:
\memoanswer{
惯性是物体的固有属性，质量是决定物体惯性大小的唯一因素，选项 B 错误；伽利略最早通过理想斜面实验得出力不是维持物体运动的必然结果，选项 C 错误。
}

\item
%% number: 15
%% paperName: 2012年山东理综第 15 题 6 分
%% typeId: 1
%% type: 单选题
%% chapter: 曲线运动
%% point: 人造地球卫星
%% method: 无
%% score: 6
%% degree: 650
%% duplicateId: 0
%% body:
2011年11月3日，“神舟八号”飞船与“天宫一号”目标飞行器成功实施了首次交会对接。任务完成后“天宫一号”经变轨升到更高的轨道，等待与“神舟九号”交会对接。变轨前和变轨完成后“天宫一号”的运行轨道均可视为圆轨道，对应的轨道半径分别为$R_{1}$、$R_{2}$，线速度大小分别为$v_{1}$、$v_{2}$。则$\frac{v_{1}}{v_{2}}$等于\xzanswer{B}
\fourchoices[answer=B]
{$\sqrt {\frac{{R_1^3}}{{R_2^3}}} $}
{$\sqrt {\frac{{{R_2}}}{{{R_1}}}} $}
{$\frac{{R_2^2}}{{R_1^2}}$}
{$\frac{{{R_2}}}{{{R_1}}}$}
%% answer: B
%% memo:
\memoanswer{
万有引力提供向心力有 $G\frac{Mm}{R^{2}}=m\frac{v^{2}}{R}$，得 $v=\sqrt{\frac{GM}{R}}$，所以 $\frac{v_{1}}{v_{2}}=\sqrt{\frac{R_{2}}{R_{1}}}$，选项 B 正确。
}

\item
%% number: 16
%% paperName: 2012年山东理综第 16 题 5 分
%% typeId: 2
%% type: 多选题
%% chapter: 运动和力
%% point: 超重失重
%% method: 无
%% score: 5
%% degree: 650
%% duplicateId: 0
%% body:
将地面上静止的货物竖直向上吊起，货物由地面运动至最高点的过程中，$v-t$图像如图所示。以下判断正确的是\xzanswer{AC}
\onepicture[h=2.8cm, align=right]{figs/16.png}
\fourchoices[answer=AC]
{前 3 s 内货物处于超重状态}
{最后 2 s 内货物只受重力作用}
{前 3 s 内与最后 2 s 内货物的平均速度相同}
{第 3 s 末至第 5 s 末的过程中，货物的机械能守恒}
%% answer: AC
%% memo:
\memoanswer{
根据 $v-t$ 图像可知，前 3 s 内货物向上做匀加速直线运动，加速度向上，处于超重状态，选项 A 正确；最后 2 s 内货物做匀减速直线运动，加速度大小为 $a=\frac{\Delta v}{\Delta t}=3\Umsq$，受重力和拉力作用，选项 B 错误；根据匀变速直线运动平均速度公式 $\overline{v}=\frac{v_0+v_t}{2}$，前 3 s 内与最后 2 s 内货物的平均速度相同，都为 $3\Ums$，选项 C 正确；第 3 s 末至第 5 s 末的过程中，货物匀速上升，机械能不守恒，选项 D 错误。
}

\item
%% number: 17
%% paperName: 2012年山东理综第 17 题 5 分
%% typeId: 2
%% type: 多选题
%% chapter: 力物体的平衡
%% point: 动态平衡
%% method: 整体与隔离
%% score: 5
%% degree: 450
%% duplicateId: 0
%% body:
如图所示，两相同轻质硬杆OO$_{1}$、OO$_{2}$可绕其两端垂直纸面的水平轴O、O$_{1}$、O$_{2}$转动，在O点悬挂一重物$M$，将两相同木块$m$紧压在竖直挡板上，此时整个系统保持静止。$F_{f}$表示木块与挡板间摩擦力的大小，$F_{N}$表示木块与挡板间正压力的大小。若挡板间的距离稍许增大后，系统仍静止且O$_{1}$、O$_{2}$始终等高，则\xzanswer{BD}
\onepicture[h=3.0cm, align=right]{figs/17.png}
\fourchoices[answer=BD]
{$F_{f}$变小}
{$F_{f}$不变}
{$F_{N}$变小}
{$F_{N}$变大}
%% answer: BD
%% memo:
\memoanswer{
把重物 $M$ 和两木块 $m$ 看成整体受力分析可得，竖直方向合力为零，始终木块与挡板间摩擦力 $2F_{f}=2mg+Mg$，选项 A 错误、B 正确；挡板间的距离稍许增大后，对结点 O 受力分析可得，轻杆弹力增大，对木块受力分析得木块与挡板间正压力增大，选项 C 错误、D 正确。
}

\item
%% number: 18
%% paperName: 2012年山东理综第 18 题 5 分
%% typeId: 2
%% type: 多选题
%% chapter: 交流电
%% point: 变压器
%% method: 无
%% score: 5
%% degree: 450
%% duplicateId: 0
%% body:
图甲是某燃气炉点火装置的原理图。转换器将直流电压转换为图乙所示的正弦交变电压，并加在一理想变压器的原线圈上，变压器原、副线圈的匝数分别为$n_{1}$、$n_{2}$。V为交流电压表。当变压器副线圈电压的瞬时值大于$5000\UV$时，就会在钢针和金属板间引发电火花进而点燃气体。以下判断正确的是\xzanswer{BC}
\onepicture[h=3.0cm, align=right]{figs/18.png}
\fourchoices[answer=BC]
{电压表的示数等于$5\UV$}
{电压表的示数等于$\frac{5}{{\sqrt 2 }}\UV$}
{实现点火的条件是$\frac{{{n_2}}}{{{n_1}}}>1000$}
{实现点火的条件是$\frac{{{n_2}}}{{{n_1}}}<1000$}
%% answer: BC
%% memo:
\memoanswer{
电压表的示数为有效值，所以等于 $U_{1}=\frac{E_{m}}{\sqrt{2}}=\frac{5\sqrt{2}}{2}\UV$，选项 A 错误、B 正确；实现点火的条件是 $U_{2}>5000\UV$，所以 $\frac{n_{2}}{n_{1}}=\frac{U_{2}}{U_{1}}>1000$，选项 C 正确、D 错误。
}

\item
%% number: 19
%% paperName: 2012年山东理综第 19 题 5 分
%% typeId: 2
%% type: 多选题
%% chapter: 电场
%% point: 等势面
%% method: 无
%% score: 5
%% degree: 450
%% duplicateId: 0
%% body:
图中虚线为一组间距相等的同心圆，圆心处固定一带正电的点电荷。一带电粒子以一定初速度射入电场，实线为粒子仅在电场力作用下的运动轨迹，a、b、c 三点是实线与虚线的交点。则该粒子\xzanswer{CD}
\onepicture[h=3.2cm, align=right]{figs/19.png}
\fourchoices[answer=CD]
{带负电}
{在 c 点受力最大}
{在 b 点的电势能大于在 c 点的电势能}
{由 a 点到 b 点的动能变化大于由 b 点到 c 点的动能变化}
%% answer: CD
%% memo:
\memoanswer{
根据粒子运动轨迹可知，粒子带正电，选项 A 错误；根据库仑定律可知，离点电荷最近时受力最大，选项 B 错误；从 b 点到 c 点电场力做正功，电势能减小，选项 C 正确；同心圆间距相等，所以 a 点到 b 点电势差大于 b 点到 c 点的电势差，所以由 a 点到 b 点的动能变化大于由 b 点到 c 点的动能变化，选项 D 正确。
}

\item
%% number: 20
%% paperName: 2012年山东理综第 20 题 5 分
%% typeId: 2
%% type: 多选题
%% chapter: 电磁感应
%% point: 电磁感应能量分析
%% method: 无
%% score: 5
%% degree: 450
%% duplicateId: 0
%% body:
如图所示，相距为$L$的两条足够长的光滑平行金属导轨与水平面的夹角为$\theta$，上端接有定值电阻$R$，匀强磁场垂直于导轨平面，磁感应强度为$B$。将质量为$m$的导体棒由静止释放，当速度达到$v$时开始匀速运动，此时对导体棒施加一平行于导轨向下的拉力，并保持拉力的功率恒为$P$，导体棒最终以$2v$的速度匀速运动。导体棒始终与导轨垂直且接触良好，不计导轨和导体棒的电阻，重力加速度为$g$。下列选项正确的是\xzanswer{AC}
\onepicture[h=3.2cm, align=right]{figs/20.png}
\fourchoices[answer=AC]
{$P=2mgv\sin\theta$}
{$P=3mgv\sin\theta$}
{当导体棒速度达到$\frac{v}{2}$时加速度大小为$\frac{g}{2}\sin\theta$}
{在速度达到$2v$以后匀速运动的过程中，$R$上产生的焦耳热等于拉力所做的功}
%% answer: AC
%% memo:
\memoanswer{
当速度达到 $v$ 时开始匀速运动，受力分析可得 $mg\sin\theta=\frac{B^{2}l^{2}v}{R}$，导体棒最终以 $2v$ 的速度匀速运动时，拉力为 $F=mg\sin\theta$，所以拉力的功率为 $P=2mgv\sin\theta$，选项 A 正确、B 错误。当导体棒速度达到 $\frac{v}{2}$ 时安培力 $F=\frac{1}{2}mg\sin\theta$，加速度为 $a=\frac{1}{2}g\sin\theta$，选项 C 正确。在速度达到 $2v$ 以后匀速运动的过程中，根据能量守恒定律，$R$ 上产生的焦耳热等于拉力所做的功加上重力做的功，选项 D 错误。
}

\item
%% number: 21
%% paperName: 2012年山东理综第 21 题 8 分
%% typeId: 4
%% type: 实验题
%% chapter: 直线运动
%% point: 打点计时器公式
%% method: 无
%% score: 8
%% degree: 450
%% duplicateId: 0
%% body:
某同学利用图甲所示的实验装置，探究物块在水平桌面上的运动规律。物块在重物的牵引下开始运动，重物落地后，物块再运动一段距离停在桌面上（尚未到达滑轮处）。从纸带上便于测量的点开始，每5个点取1个计数点，相邻计数点间的距离如图乙所示。打点计时器电源的频率为$50\UHz$。
\onepicture[h=3.0cm, align=right]{figs/21.png}
\begin{enumerate}
	\item
	通过分析纸带数据，可判断物块在两相邻计数点\_\_\_\_\_\_\_\_和\_\_\_\_\_\_\_\_之间某时刻开始减速。
	\item
	计数点5对应的速度大小为\_\_\_\_\_\_\_\_$\Ums$，计数点6对应的速度大小为\_\_\_\_\_\_\_\_$\Ums$。（保留三位有效数字）
	\item
	物块减速运动过程中加速度的大小为$a=$\_\_\_\_\_\_\_\_$\Umsq$，若用$a$来计算物块与桌面间的动摩擦因数（$g$为重力加速度），则计算结果比动摩擦因数的真实值\_\_\_\_\_\_\_\_（填“偏大”或“偏小”）。
\end{enumerate}
%% answer: 
\jdanswer{
\begin{enumerate}
	\item
	6；7（或7；6）

	\item
	1.00；1.20

	\item
	2.00；偏大

\end{enumerate}
}
%% memo:
\memoanswer{
①由于计数点前后的间隔距离都小于它们的间隔距离，说明计数点 6 之前物块在加速，计数点 7 之后物块在减速，则开始减速的时刻在 6 和 7 之间。

②计数点 5 对应的速度等于 4 和 6 间的平均速度 $v_{5}=\frac{(9.00+11.01)\times10^{-2}}{0.2}\Ums=1.00\Ums$；同理 $v_{4}=\frac{(7.01+9.00)\times10^{-2}}{0.2}\Ums=0.80\Ums$；又 $v_{5}=\frac{v_{4}+v_{6}}{2}$ 可解得计数点 6 对应的速度大小为 $1.20\Ums$。

③在减速阶段 $\Delta x=2.00\Ucm$，则加速度为 $a=\frac{\Delta x}{T^{2}}=\frac{2.00\times10^{-2}}{0.1^{2}}\Umsq=2.00\Umsq$。在减速阶段产生加速度的力是滑动摩擦力和纸带受到的阻力，所以计算结果比动摩擦因数的真实值“偏大”。
}

\item
%% number: 21
%% paperName: 2012年山东理综第 21 题 5 分
%% typeId: 4
%% type: 实验题
%% chapter: 电路
%% point: 伏安法测电阻
%% method: 无
%% score: 5
%% degree: 450
%% duplicateId: 0
%% body:
在测量金属丝电阻率的实验中，可供选用的器材如下：

待测金属丝：$R_{x}$（阻值约$4\UO$，额定电流约$0.5\UA$）；

电压表：V（量程$3\UV$，内阻约$3\UkO$）；

电流表：$A_{1}$（量程$0.6\UA$，内阻约$0.2\UO$）；

$A_{2}$（量程$3\UA$，内阻约$0.05\UO$）；

电源：$E_{1}$（电动势$3\UV$，内阻不计）；

$E_{2}$（电动势$12\UV$，内阻不计）；

滑动变阻器：$R$（最大阻值约$20\UO$）；

螺旋测微器；毫米刻度尺；开关 S；导线。
\onepicture[h=2.8cm, align=right]{figs/21b.png}
\begin{enumerate}
	\item
	用螺旋测微器测量金属丝的直径，示数如图所示，读数为\_\_\_\_\_\_\_\_$\Umm$。
	\item
	若滑动变阻器采用限流接法，为使测量尽量精确，电流表应选\_\_\_\_\_\_\_\_、电源应选\_\_\_\_\_\_\_\_（均填器材代号），在虚线框内完成电路原理图。
\end{enumerate}
%% answer: 
\jdanswer{
\begin{enumerate}
	\item
	$1.773\Umm$（1.771$\sim$1.775 均可）

	\item
	电流表应选 $A_{1}$，电源应选 $E_{1}$；电路采用电流表外接、滑动变阻器限流接法。

\end{enumerate}
}
%% memo:
\memoanswer{
①螺旋测微器的固定刻度为 $1.5\Umm$，可动刻度为 $27.3\times0.01\Umm=0.273\Umm$，读数为 $1.773\Umm$（1.771$\sim$1.775 均可）。

②电压表量程 $3\UV$，所以电源应选 $E_{1}$，通过待测金属丝的最大电流约为 $I=\frac{3}{4}\UA=0.75\UA$，所以电流表应选 $A_{1}$。电路原理图如图所示。

\onepicture[h=3.6cm, align=right]{figs/21b_answer.png}
}

\item
%% number: 22
%% paperName: 2012年山东理综第 22 题 15 分
%% typeId: 6
%% type: 计算题
%% chapter: 功和能
%% point: 动能定理
%% method: 无
%% score: 15
%% degree: 450
%% duplicateId: 0
%% body:
如图所示，一工件置于水平地面上，其AB段为一半径$R=1.0\Um$的光滑圆弧轨道，BC段为一长度$L=0.5\Um$的粗糙水平轨道，二者相切于B点，整个轨道位于同一竖直平面内，P点为圆弧轨道上的一个确定点。一可视为质点的物块，其质量$m=0.2\Ukg$，与$BC$间的动摩擦因数$\mu_{1}=0.4$。工件质量$M=0.8\Ukg$，与地面间的动摩擦因数$\mu_{2}=0.1$。（取$g=10\Umsq$）
\begin{enumerate}
	\item
	若工件固定，将物块由P点无初速度释放，滑至$C$点时恰好静止，求P、C两点间的高度差$h$。
	\item
	若将一水平恒力$F$作用于工件，使物块在$P$点与工件保持相对静止，一起向左做匀加速直线运动。①求$F$的大小。②当速度$v=5\Ums$时，使工件立刻停止运动（即不考虑减速的时间和位移），物块飞离圆弧轨道落至BC段，求物块的落点与B点间的距离。
\end{enumerate}
\onepicture[h=3.6cm, align=right]{figs/22.png}
%% answer: 
\jdanswer{
\begin{enumerate}
	\item
	$0.2\Um$

	\item
	①$8.5\UN$

	②$0.4\Um$

\end{enumerate}
}
%% memo:
\memoanswer{
（1）物块从 P 点下滑经 B 点至 C 点的整个过程，根据动能定理得

$mgh-\mu_{1}mgL=0$ ①

代入数据得

$h=0.2\Um$ ②

（2）①设物块的加速度大小为 $a$，P 点与圆心的连线与竖直方向间的夹角为 $\theta$，由几何关系可得

$\cos\theta=\frac{R-h}{R}$ ③

根据牛顿第二定律，对物体有

$mg\tan\theta=ma$ ④

对工件和物体整体有

$F-\mu_{2}(M+m)g=(M+m)a$ ⑤

联立②③④⑤式，代入数据得 $F=8.5\UN$ ⑥

②设物体平抛运动的时间为 $t$，水平位移为 $x_{1}$，物块落点与 B 间的距离为 $x_{2}$，由运动学公式可得

$h=\frac{1}{2}gt^{2}$ ⑦

$x_{1}=vt$ ⑧

$x_{2}=x_{1}-R\sin\theta$ ⑨

联立②③⑦⑧⑨式，代入数据得 $x_{2}=0.4\Um$ ⑩
}

\item
%% number: 23
%% paperName: 2012年山东理综第 23 题 18 分
%% typeId: 6
%% type: 计算题
%% chapter: 磁场
%% point: 带电粒子在组合场中的运动
%% method: 无
%% score: 18
%% degree: 450
%% duplicateId: 0
%% body:
如图甲所示，相隔一定距离的竖直边界两侧为相同的匀强磁场区，磁场方向垂直纸面向里，在边界上固定两长为$L$的平行金属极板MN和PQ，两极板中心各有一小孔S$_{1}$、S$_{2}$，两极板间电压的变化规律如图乙所示，正反向电压的大小均为$U_{0}$，周期为$T_{0}$。在$t=0$时刻将一个质量为$m$、电荷量为$-q$（$q>0$）的粒子由S$_{1}$静止释放，粒子在电场力的作用下向右运动，在$t=\frac{T_{0}}{2}$时刻通过S$_{2}$垂直于边界进入右侧磁场区。（不计粒子重力，不考虑极板外的电场）
\begin{enumerate}
	\item
	求粒子到达S$_{2}$时的速度大小$v$和极板间距$d$。
	\item
	为使粒子不与极板相撞，求磁感应强度的大小应满足的条件。
	\item
	若已保证了粒子未与极板相撞，为使粒子在$t=3T_{0}$时刻再次到达S$_{2}$，且速度恰好为零，求该过程中粒子在磁场内运动的时间和磁感应强度的大小。
\end{enumerate}
\onepicture[h=3.4cm, align=right]{figs/23.png}
%% answer: 
\jdanswer{
\begin{enumerate}
	\item
	$v=\sqrt{\frac{2qU_{0}}{m}}$，$d=\frac{T_{0}}{4}\sqrt {\frac{{2q{U_0}}}{m}} $

	\item
	$B<\frac{4}{L}\sqrt {\frac{{2m{U_0}}}{q}} $

	\item
	$t=\frac{7T_{0}}{4}$，$B=\frac{{8\pi m}}{{7q{T_0}}}$

\end{enumerate}
}
%% memo:
\memoanswer{
（1）粒子由 $S_{1}$ 至 $S_{2}$ 的过程中，根据动能定理得

$qU_{0}=\frac{1}{2}mv^{2}$ ①

由①式得

$v=\sqrt{\frac{2qU_{0}}{m}}$ ②

设粒子的加速度大小为 $a$，由牛顿第二定律得

$q\frac{U_{0}}{d}=ma$ ③

由运动学公式得 $d=\frac{1}{2}a\left(\frac{T_{0}}{2}\right)^{2}$ ④

联立③④式得

$d=\frac{T_{0}}{4}\sqrt{\frac{2qU_{0}}{m}}$ ⑤

（2）设磁感应强度大小为 $B$，粒子在磁场中做匀速圆周运动的半径为 $R$，由牛顿第二定律得

$qvB=m\frac{v^{2}}{R}$ ⑥

要使粒子在磁场中运动时不与极板相撞，须满足

$2R>\frac{L}{2}$ ⑦

联立②⑥⑦式得 $B<\frac{4}{L}\sqrt{\frac{2mU_{0}}{q}}$ ⑧

（3）设粒子在两边界之间无场区向左匀速运动的过程用时为 $t_{1}$，有

$d=vt_{1}$ ⑨

联立②⑤⑨式得 $t_{1}=\frac{T_{0}}{4}$ ⑩

若粒子再次到达 $S_{2}$ 时速度恰好为零，粒子回到极板间应做匀减速运动，设匀减速运动的时间为 $t_{2}$，根据运动学公式得 $d=\frac{v}{2}t_{2}$ (11)

联立⑨⑩(11)式得 $t_{2}=\frac{T_{0}}{2}$ (12)

设粒子在磁场中运动的时间为 $t$

$t=3T_{0}-\frac{T_{0}}{2}-t_{1}-t_{2}$ (13)

联立⑩(12)(13)式得 $t=\frac{7T_{0}}{4}$ (14)

设粒子在匀强磁场中做匀速圆周运动的周期为 $T$，由⑥式结合运动学公式得

$T=\frac{2\pi m}{qB}$ (15)

由题意得 $T=t$ (16)

联立(14)(15)(16)式得 $B=\frac{8\pi m}{7qT_{0}}$ (17)
}

\item
%% number: 36
%% paperName: 2012年山东理综第 36 题 4 分
%% typeId: 2
%% type: 多选题
%% chapter: 气体性质
%% point: 分子动理论
%% method: 无
%% score: 4
%% degree: 650
%% duplicateId: 0
%% body:
以下说法正确的是\xzanswer{AB}
\fourchoices[answer=AB]
{水的饱和汽压随温度的升高而增大}
{扩散现象表明，分子在永不停息地运动}
{当分子间距离增大时，分子间引力增大，分子间斥力减小}
{一定质量的理想气体，在等压膨胀过程中，气体分子的平均动能减小}
%% answer: AB
%% memo:
\memoanswer{
当分子间距离增大时，分子间引力减小，分子间斥力也减小，选项 C 错误；一定质量的理想气体，在等压膨胀过程中，温度升高，气体分子的平均动能增大，选项 D 错误。
}

\item
%% number: 36
%% paperName: 2012年山东理综第 36 题 4 分
%% typeId: 3
%% type: 填空题
%% chapter: 气体性质
%% point: 玻意耳定律
%% method: 无
%% score: 4
%% degree: 650
%% duplicateId: 0
%% body:
如图所示，粗细均匀、导热良好、装有适量水银的U形管竖直放置，右端与大气相通，左端封闭气柱长$l_{1}=20\Ucm$（可视为理想气体），两管中水银面等高。现将右端与一低压舱（未画出）接通，稳定后右管水银面高出左管水银面$h=10\Ucm$。（环境温度不变，大气压强$p_{0}=75\UcmHg$）
\begin{enumerate}
	\item
	求稳定后低压舱内的压强（用“cmHg”作单位）。
	\item
	此过程左管内的气体对外界\_\_\_\_\_\_\_\_（填“做正功”“做负功”或“不做功”），气体将\_\_\_\_\_\_\_\_（填“吸热”或“放热”）。
\end{enumerate}
\onepicture[h=3.4cm, align=right]{figs/36.png}
%% answer: 
\jdanswer{
\begin{enumerate}
	\item
	$50\UcmHg$

	\item
	做正功；吸热

\end{enumerate}
}
%% memo:
\memoanswer{
（1）设 U 形管横截面积为 $S$，右端与大气相通时左管中封闭气体压强为 $P_{1}$，右端与一低压舱接通后左管中封闭气体压强为 $P_{2}$，气柱长度为 $l_{2}$，稳定后低压舱内的压强为 $P$。左管中封闭气体发生等温变化，根据玻意耳定律得：$p_{1}V_{1}=p_{2}V_{2}$ ①

$P_{1}=P_{0}$ ②

$P_{2}=P+P_{h}$ ③

$V_{1}=l_{1}S$ ④

$V_{2}=l_{2}S$ ⑤

由几何关系得 $h=2(l_{2}-l_{1})$ ⑥

联立①②③④⑤⑥式，代入数据得

$P=50\UcmHg$ ⑦

（2）由于左管内的气体体积增大，对外界做正功；而温度不变，由热力学第一定律 $\Delta U=W+Q$，气体将吸热。
}

\item
%% number: 37
%% paperName: 2012年山东理综第 37 题 4 分
%% typeId: 6
%% type: 计算题
%% chapter: 机械波
%% point: 波的图象
%% method: 无
%% score: 4
%% degree: 650
%% duplicateId: 0
%% body:
一列简谐横波沿$x$轴正方向传播，$t=0$时刻的波形如图所示，介质中质点$P$、$Q$分别位于$x=2\Um$、$x=4\Um$处。从$t=0$时刻开始计时，当$t=15\Us$时质点$Q$刚好第4次到达波峰。
\begin{enumerate}
	\item
	求波速。
	\item
	写出质点$P$做简谐运动的表达式（不要求推导过程）。
\end{enumerate}
\onepicture[h=3.0cm, align=right]{figs/37.png}
%% answer: 
\jdanswer{
\begin{enumerate}
	\item
	$1\Ums$

	\item
	$y=0.2\sin(0.5\pi t)\Um$

\end{enumerate}
}
%% memo:
\memoanswer{
①设简谐横波的波速为 $v$，波长为 $\lambda$，周期为 $T$，由图像知，$\lambda=4\Um$。由题意得

$t=3T+\frac{3}{4}T$ ①

$v=\frac{\lambda}{T}$ ②

联立①②式，代入数据得 $v=1\Ums$ ③

②质点 P 做简谐运动的表达式为

$y=0.2\sin(0.5\pi t)\Um$ ④
}

\item
%% number: 37
%% paperName: 2012年山东理综第 37 题 4 分
%% typeId: 6
%% type: 计算题
%% chapter: 光学
%% point: 光的折射
%% method: 无
%% score: 4
%% degree: 650
%% duplicateId: 0
%% body:
如图所示，一玻璃球体的半径为$R$，O为球心，AB为直径。来自B点的光线BM在M点射出，出射光线平行于AB，另一光线BN恰好在N点发生全反射。已知$\angle$ABM$=30^{\circ}$，求：
\begin{enumerate}
	\item
	玻璃的折射率。
	\item
	球心 O 到 BN 的距离。
\end{enumerate}
\onepicture[h=3.2cm, align=right]{figs/37b.png}
%% answer: 
\jdanswer{
\begin{enumerate}
	\item
	$\sqrt 3 $

	\item
	$\frac{{\sqrt 3 }}{3}R$

\end{enumerate}
}
%% memo:
\memoanswer{
①设光线 BM 在 M 点的入射角为 $i$，折射角为 $r$，由几何关系可知，$i=30^{\circ}$，$r=60^{\circ}$。根据折射定律得

$n=\frac{\sin r}{\sin i}$ ①

代入数据得

$n=\sqrt{3}$ ②

②光线 BN 恰好在 N 点发生全反射，则 $\angle BNO$ 为临界角 $C$。

$\sin C=\frac{1}{n}$ ③

设球心到 BN 的距离为 $d$，由几何关系可知

$d=R\sin C$ ④

联立③④式得 $d=\frac{\sqrt{3}}{3}R$ ⑤
}

\item
%% number: 38
%% paperName: 2012年山东理综第 38 题 2 分
%% typeId: 3
%% type: 填空题
%% chapter: 原子和原子核
%% point: 玻尔模型
%% method: 无
%% score: 2
%% degree: 650
%% duplicateId: 0
%% body:
氢原子第$n$能级的能量为$E_{n}=\frac{{{E_1}}}{{{n^2}}}$，其中$E_{1}$为基态能量。当氢原子由第4能级跃迁到第2能级时，发出光子的频率为$\nu_{1}$；若氢原子由第2能级跃迁到基态，发出光子的频率为$\nu_{2}$，则$\frac{{{\nu_1}}}{{{\nu_2}}}=$\_\_\_\_\_\_\_。
%% answer: 
\jdanswer{
$\frac{1}{4}$
}
%% memo:
\memoanswer{
根据跃迁公式 $h\nu=E_{m}-E_{n}$ 即可解得。
}

\item
%% number: 38
%% paperName: 2012年山东理综第 38 题 6 分
%% typeId: 6
%% type: 计算题
%% chapter: 运动和力
%% point: 动量守恒定律
%% method: 无
%% score: 6
%% degree: 450
%% duplicateId: 0
%% body:
光滑水平轨道上有三个木块A、B、C，质量分别为$m_{A}=3m$、$m_{B}=m_{C}=m$，开始时B、C均静止，A以初速度向右运动，A与B相撞后分开，B又与C发生碰撞并粘在一起，此后A与B间的距离保持不变。求B与C碰撞前B的速度大小。
\onepicture[h=1.6cm, align=right]{figs/38.png}
%% answer: 
\jdanswer{
$\frac{6}{5}v_{0}$
}
%% memo:
\memoanswer{
设 AB 碰撞后 A 的速度为 $v_{A}$，B 与 C 碰撞前 B 的速度为 $v_{B}$，B 与 C 碰撞后粘在一起的速度为 $v$，由动量守恒定律得

对 A、B 木块：$m_{A}v_{0}=m_{A}v_{A}+m_{B}v_{B}$ ①

对 B、C 木块：$m_{B}v_{B}=(m_{B}+m_{C})v$ ②

由 A 与 B 间的距离保持不变可知

$v_{A}=v$ ③

联立①②③式，代入数据得

$v_{B}=\frac{6}{5}v_{0}$ ④
}

\end{enumerate}

\end{document}
